\subsection{Ersatzmodell (Surrogat) aus den Strömungsrechnungen und die Frage der Stufung}\label{sec:tauschersurrogat}

\paragraph{Warum ein Surrogat?} Die Auslegung des Vorwärmers braucht für jeden Entwurf Nu und $f$ auf beiden
Seiten. Dafür gab es bisher drei Wege, und alle haben einen Haken:
\begin{enumerate}
\item \textbf{Korrelationen aus der Literatur} sind schnell, streuen aber stark. Bei Re 3000 liefern Renon, Li 2020
  und Li 2022 für denselben Gyroid $f$ = 0,54, 0,15 und 0,81, also mehr als den Faktor 5, und Nu = 75, 63 und 79
  \cite{renon2024,alsafadi2026}. Sie gelten außerdem nur für die Zellen und Wände, mit denen sie gemessen oder
  gerechnet wurden.
\item \textbf{Eigene Strömungsrechnungen} passen genau zur eigenen Geometrie, kosten aber 1--70 Minuten je Punkt.
  Das Segmentmodell weiter unten wertet je Variante etwa 440\,000 Kombinationen aus Annäherung, Geschwindigkeit,
  Durchgängen, Zelle und Segment aus. Mit je zwei Minuten Rechenzeit wären das fast zwei Jahre.
\item \textbf{Eine eigene Korrelation} (Potenzgesetz $\mathrm{Nu} = c\,\mathrm{Re}^m\Pr^n$) an die Rechenpunkte
  anzupassen wäre schnell, legt die Form aber vorher fest und sagt nicht, wie sicher sie zwischen den Punkten ist.
\end{enumerate}
Ein \emph{Surrogat} (Ersatzmodell) verbindet die Vorteile. Es wird einmal an wenigen teuren Rechnungen angelernt und
liefert danach in Mikrosekunden Werte für beliebige Eingänge, samt einer Angabe, wie verlässlich diese sind. Hier
geht es um eine konkrete Frage: Lohnt sich eine Stufung der Zellen längs des Wegs? Die Literatur gibt dafür nur einen
pauschalen Faktor (Nu $\times$ 1,28 \cite{oh2025}). Um das selbst zu prüfen, muss man wissen, wie Nu und $f$ von der
Zellgröße abhängen, und genau das lernt das Surrogat aus den eigenen Rechnungen.

\paragraph{Wie ein Gauß-Prozess funktioniert.} Als Surrogat dient ein Gauß-Prozess (Kriging). Er nimmt nicht eine
Formel an, sondern nur, dass ähnliche Eingänge ähnliche Ergebnisse geben. Diese Ähnlichkeit beschreibt ein Kern
\[
  k(x,x') = \sigma^2 \exp\Big(-\sum_i \frac{(x_i-x_i')^2}{2\ell_i^2}\Big) + \sigma_n^2\,\delta_{xx'}.
\]
Die Längenskala $\ell_i$ sagt, über welche Entfernung in Eingang $i$ sich das Ergebnis merklich ändert. Ein großes
$\ell_i$ heißt, Eingang $i$ ist fast bedeutungslos. $\sigma_n^2$ fängt das Rauschen der Rechnungen auf, etwa die
zeitliche Schwankung der Mischtemperatur. Alle Größen werden aus den Daten bestimmt (größte Randwahrscheinlichkeit).
An einem neuen Punkt liefert der Prozess einen Mittelwert, eine gewichtete Mischung der bekannten Punkte, und eine
Standardabweichung. Diese ist nahe an Rechenpunkten klein und wächst in Lücken. Damit sieht man, wo das Modell rät.
Die Längenskalen sind selbst ein Ergebnis: Sie zeigen, welcher Eingang wirklich zählt.

\paragraph{Warum mehrstufig?} Feine Gitter sind genau, aber teuer (n144: 30--70 Minuten, 8\,GB Grafikspeicher); grobe
Gitter sind billig (n48: 1--2 Minuten), aber systematisch zu hoch. Ein mehrstufiges Surrogat (Co-Kriging
\cite{kennedy2000}) nutzt beides. Viele grobe Rechnungen beschreiben die \emph{Form} der Abhängigkeit von Re, Pr und
Zelle. Wenige feine Rechnungen bestimmen die \emph{Verschiebung} durch das Gitter. Das geht gut, wenn der Gitterfehler
sich langsam und glatt über den Eingangsraum ändert, und das wird unten geprüft.

\paragraph{Versuchsplan.} Für Gyroid und Diamond wurden je 14 Punkte per Latin Hypercube \cite{mckay1979} gewählt:
Re 500--10\,000 (logarithmisch), Pr 1,2--2,0, Zelle $a$ = 0,75--1,6 Standardzelle, Wand fest 0,5\,mm. Diese Punkte
liefen auf dem groben Gitter (48 Zellen je $d_h$). Dazu kamen drei Kontrollpunkte auf dem mittleren Gitter (96) und je
ein Lauf auf dem feinen Gitter (144 Zellen je $d_h$, 62,5 Mio.\ Zellen). Mit den Läufen aus
Abschnitt~\ref{sec:tauschercfd} sind es 40 Läufe. Die Reihe lief unbeaufsichtigt über eine Auftragsliste, die
jederzeit angehalten werden konnte. Ein zweiter KI-Agent prüfte Speicher, Platte und Divergenz. Es gab keinen
Fehlschlag.

\paragraph{Aufbau.} Damit der Gauß-Prozess nur die Abweichung lernen muss, wird auf die Korrelation von Renon
\cite{renon2024} bezogen:
\[
  y_\mathrm{Nu} = \ln\frac{\mathrm{Nu}}{\mathrm{Nu}_\mathrm{Renon}(\mathrm{Re},\Pr)},\qquad
  y_f = \ln\frac{f}{f_\mathrm{Renon}(\mathrm{Re})},\qquad x = (\ln\mathrm{Re},\ \Pr,\ a).
\]
\begin{enumerate}
\item \textbf{Grobe Stufe:} je Bauart ein Gauß-Prozess mit anisotropem RBF-Kern und Rauschterm über die 15 Punkte auf
  48 Zellen. Die Kreuzprüfung (jeweils ein Punkt weggelassen) ergibt für Nu 2,8\,\% (Gyroid) bzw.\ 0,9\,\% (Diamond)
  mittleren Fehler. Für $f$ sind es 9,5\,\% bzw.\ 1,1\,\%; die Gyroid-Reibung streut also deutlich.
\item \textbf{Gitterkorrektur:} An jedem Kontrollpunkt ist $d_{96} = y_{96} - y_{48}(x)$. Über Re verläuft das als
  Gerade in $\ln\mathrm{Re}$ (Rest 0,03 in $\ln$). Das entspricht einem Co-Kriging nach Kennedy und O'Hagan
  \cite{kennedy2000} mit $\rho = 1$ in der logarithmischen Größe.
\item \textbf{Feinste Stufe:} Die Reihe 48/96/144 bei Re 3000 ist \emph{nicht} im asymptotischen Bereich. Es gilt
  $d_{96}/d_{144} = 0{,}58$, ein Gitterfehler $\propto n^{-p}$ verlangt aber mehr als $\ln 2/\ln 3 = 0{,}63$; keine
  positive Ordnung passt. Darum gibt es zwei obere Stufen. \emph{n144} ist gerechnet: $d_{96}(x)$ wird mit
  $d_{144}/d_{96} \approx 1{,}7$ skaliert. \emph{Gitter $\to\infty$} ist geschätzt, mit Richardson nur über den letzten
  Schritt und der Ordnung aus dem Plattentest:
  $d_\infty = d_{144} + (d_{144}-d_{96})/(1{,}5^{1{,}33}-1)$. Als Unsicherheit dient der halbe Abstand zwischen beiden
  Stufen.
\end{enumerate}

\begin{table}[ht]\centering\small
\begin{tabular}{llcccc}\toprule
& & \multicolumn{2}{c}{$\mathrm{Nu}/\mathrm{Nu}_\mathrm{Renon}$} & \multicolumn{2}{c}{$f/f_\mathrm{Renon}$} \\
Bauart & Re & n144 & Gitter $\to\infty$ & n144 & Gitter $\to\infty$ \\\midrule
Gyroid & 1000 & 1,20 & $1{,}03\pm11\,\%$ & $0{,}94\pm15\,\%$ & (unsicher) \\
Gyroid & 3000 & 1,29 & $1{,}03\pm14\,\%$ & $1{,}06\pm15\,\%$ & (unsicher) \\
Gyroid & 8000 & 1,49 & $1{,}14\pm16\,\%$ & $1{,}35\pm15\,\%$ & (unsicher) \\
Diamond & 1000 & 1,14 & $0{,}97\pm12\,\%$ & $0{,}98\pm6\,\%$ & (unsicher) \\
Diamond & 3000 & 1,23 & $1{,}00\pm14\,\%$ & $1{,}01\pm6\,\%$ & (unsicher) \\
Diamond & 8000 & 1,40 & $1{,}08\pm16\,\%$ & $1{,}06\pm6\,\%$ & (unsicher) \\\bottomrule
\end{tabular}
\caption{Surrogat bei Pr 1,5 und Standardzelle. Die extrapolierte Reibung ist nicht verwertbar: Sie fiele unter Renon,
weil $f$ zwischen den Gittern nicht monoton konvergiert.}
\end{table}

\bild{tauscher_surrogat.png}{1.0}{Surrogat für Gyroid (blau) und Diamond (orange). Links Nu: punktiert grobe Stufe,
gestrichelt n144, durchgezogen mit Band Gitter $\to\infty$, grau gestrichelt Renon. Mitte $f$ (n144 mit Band).
Rechts die Kreuzprüfung der groben Stufe.}
\lesehilfe{Ein \emph{Surrogat} ist ein schnelles Ersatzmodell. Es lernt aus 40 aufwendigen
Strömungsrechnungen und liefert dann für beliebige Werte sofort eine Schätzung samt Unsicherheit. Links die
Nusselt-Zahl (Wärmeübergang) über der \emph{Reynolds-Zahl}, einem Maß dafür, wie schnell und verwirbelt die
Strömung ist; beide Achsen sind logarithmisch. Die Punkte sind Rechnungen (Kreise grob, Quadrate fein, Rauten am
feinsten), die Linien das Surrogat, das helle Band seine Unsicherheit, die grau gestrichelte Linie die Literatur.
In der Mitte dasselbe für den Reibungsbeiwert $f$, der den Druckverlust bestimmt. Rechts die Kreuzprüfung: Jede
Rechnung wurde einmal weggelassen und aus den übrigen vorhergesagt. Liegen die Punkte auf der Diagonale, stimmt die
Vorhersage; die Abweichung beträgt nur 1 bis 3\,\%.}


\paragraph{Was das Surrogat sagt.}
\begin{itemize}
\item Auf das unendlich feine Gitter extrapoliert, landet Nu bei der Korrelation von Renon (0,97--1,14). Der
  Überschuss von 1,5 bei 96 Zellen war also Gitterfehler, nicht Physik. Die Schätzung von 1,2 aus nur zwei Stufen war
  zu hoch.
\item Bei gleichem $d_h$ hängen Nu und $f$ kaum von der Zellgröße ab: Die Längenskala des Kerns für $a$ liegt beim
  Gyroid am oberen Rand. Wer die Zelle verkleinert, verkleinert $d_h$, ändert die Kennzahlen aber kaum.
\item Gyroid und Diamond sind beim Wärmeübergang gleichauf. Beim Druckverlust liegt der Gyroid bei hoher Re höher.
\end{itemize}

\paragraph{Anwendung: lohnt sich die Stufung?} Dazu wurde der Vorwärmer aus Abschnitt~\ref{sec:vorwaermer}
als Segmentmodell neu aufgesetzt (\texttt{stellarator/vorwaermer\_stufung.py}):
\begin{itemize}
\item Gegenstrom, Zwei-Fluid-TPMS, beide Kanalnetze teilen die Zelle (gleiches $d_h$ beiderseits).
\item 8 Segmente gleicher Wärmeleistung, je Segment lokale Temperaturen, Stoffwerte, Pr und LMTD.
\item Nu und $f$ aus dem Surrogat (Stufe n144; Nu zusätzlich Gitter $\to\infty$ als Band).
\item \emph{Gleichmäßig:} eine Zelle für alle Segmente. \emph{Abgestuft:} jedes Segment wählt seine Zelle selbst, und
  zwar so, dass $P_i + \lambda V_i$ minimal ist, mit $\lambda$ abgetastet.
\item Frei sind außerdem Annäherung, Kondensatgeschwindigkeit und Zahl der Durchgänge der warmen Seite.
\end{itemize}

\bild{vorwaermer_stufung.png}{0.85}{Netto-Mehrleistung über dem Bauvolumen. Abgestuft (gestrichelt) und gleichmäßig
fallen zusammen. Zum Vergleich Renon und Renon mit Nu $\times$ 1,28 (Stufung nach Oh et al.) im selben Modell.}
\lesehilfe{Aufbau wie bei der Netto-Mehrleistung weiter oben, jetzt mit dem Surrogat. Die
gestrichelten Linien gehören zu Tauschern, deren Zellgröße sich entlang des Wegs in acht Abschnitten ändert
(\emph{abgestuft}). Die durchgezogenen gehören zu Tauschern mit überall gleichen Zellen. Sie liegen praktisch
übereinander. Das Band reicht vom fein gerechneten bis zum hochgerechneten Wärmeübergang. Die Stufung bringt also
keinen Vorteil; die kleinste Zelle ist überall die beste.}


\paragraph{Ergebnis.}
\begin{itemize}
\item \textbf{Die Stufung bringt nichts.} In jedem Segment gewinnt die kleinste zulässige Zelle ($a = 0{,}75\,a_\mathrm{std}$,
  $d_h = 1{,}33$\,mm). Damit ist der abgestufte Entwurf identisch mit dem gleichmäßigen. Der Grund: Die Pumpleistung
  ist hier klein (unter 1\,MW). Verloren geht die Leistung vor allem über die Annäherung am warmen Ende, und die
  verringert man am billigsten mit mehr Fläche je Volumen, also mit feinen Zellen überall.
\item \textbf{Bauvolumen:} 95\,\% der Obergrenze von 73,3\,MW erreicht der Vorwärmer mit 1,3\,m$^3$ (n144) bis
  1,5\,m$^3$ (Gitter $\to\infty$). Bei 1\,m$^3$ liefern CFD und Renon 66,0--67,6\,MW, der Unterschied ist klein.
\item Der Literaturfaktor für die Stufung (Nu $\times$ 1,28 \cite{oh2025}) entspricht hier fast genau dem Abstand
  zwischen Renon und dem n144-Surrogat. Die von Oh et al.\ berichtete Verbesserung lässt sich auf Zellebene nicht
  nachbilden. Sie muss aus Effekten stammen, die dieses Modell nicht enthält, etwa der Verteilung in den Sammlern.
\item \textbf{Grenzen:} Die beste Zelle liegt am Rand des Surrogatbereichs; noch feinere Zellen sind ungerechnet. Ob
  die Wand von 0,5\,mm die Druckdifferenz von 3\,MPa trägt, ist nicht geprüft. Die Fertigung so feiner Zellen ist offen.
  Eine Stufung kann sich dort lohnen, wo die Pumpleistung zählt, zum Beispiel auf der Dampfseite des Kondensators.
\end{itemize}
